\documentclass[11pt,a4paper]{article}
\usepackage[italian]{babel}
\usepackage[margin=2.5cm]{geometry}
\usepackage{parskip}
\usepackage{amsmath, amssymb, amsthm}
\usepackage{booktabs}
\usepackage{array}
\usepackage{xcolor}
\usepackage{needspace}
\usepackage{enumitem}
\usepackage{listings}
\usepackage{tikz}
\usepackage[most]{tcolorbox}
\usepackage[hidelinks]{hyperref}
\usetikzlibrary{decorations.pathreplacing, shapes.geometric, arrows.meta, positioning}

% Niente vedove/orfane: i paragrafi non lasciano righe isolate a fine/inizio pagina
\widowpenalty=10000
\clubpenalty=10000

% Elenchi più compatti
\setlist{itemsep=2pt, parsep=0pt, topsep=4pt, partopsep=0pt}

% --- Riquadri con bordo nero, senza numero (si spezzano tra due pagine) ---
% Uso: \begin{definizione} ... \end{definizione}
%      \begin{teorema}[Nome] ... \end{teorema}  (il nome è facoltativo)
\tcbset{nota/.style={enhanced jigsaw, breakable, colback=white,
  colframe=black, boxrule=0.5pt, sharp corners}}
\NewTColorBox{definizione}{}{nota,
  before upper={\textbf{Definizione.}\ }}
\NewTColorBox{teorema}{o}{nota,
  before upper={\textbf{Teorema\IfValueT{#1}{ (#1)}.}\ }}
\NewTColorBox{proposizione}{o}{nota,
  before upper={\textbf{Proposizione\IfValueT{#1}{ (#1)}.}\ }}
\NewTColorBox{proprieta}{o}{nota,
  before upper={\textbf{Proprietà\IfValueT{#1}{ (#1)}.}\ }}
\NewTColorBox{assioma}{o}{nota, boxrule=1pt,
  before upper={\textbf{Assioma\IfValueT{#1}{ (#1)}.}\ }}

% --- Ambienti semplici (senza riquadro) ---
% Il titolo sta in cima al blocco, anche se il contenuto inizia con un riquadro o
% con codice; e non resta mai da solo in fondo a una pagina.
% Uso: \begin{esempio}[Titolo facoltativo] ... \end{esempio}
\newcommand{\newrunin}[2]{%
  \NewDocumentEnvironment{#1}{o}{%
    \par\Needspace{4\baselineskip}\addvspace{\medskipamount}\noindent
    \textbf{#2}\IfValueT{##1}{ (##1)}\textbf{.}\ \ignorespaces}%
  {\par\addvspace{\medskipamount}}}
\newrunin{esempio}{Esempio}
\newrunin{esercizio}{Esercizio}
\newrunin{osservazione}{Osservazione}

% V e F nelle tavole di verità
\newcommand{\V}{\textsf{V}}
\newcommand{\F}{\textsf{F}}

% --- Codice C ---
% Uso: \begin{codice} ... \end{codice}      codice C numerato
%      \begin{comando} ... \end{comando}    comandi da terminale
%      \begin{risultato} ... \end{risultato} output di un programma
%      \lstinline|printf("ciao");|           codice dentro al testo
\lstdefinestyle{c}{
  language=C,
  basicstyle=\ttfamily\small,
  keywordstyle=\color{blue}\bfseries,
  commentstyle=\color{gray}\itshape,
  stringstyle=\color{red!70!black},
  numbers=left,
  numberstyle=\tiny\color{gray},
  numbersep=8pt,
  columns=flexible,
  keepspaces=true,
  breaklines=true,
  showstringspaces=false,
  tabsize=4,
  morekeywords={include, define},
  % lettere accentate nei commenti
  literate={à}{{\`a}}1 {è}{{\`e}}1 {é}{{\'e}}1 {ì}{{\`i}}1
           {ò}{{\`o}}1 {ù}{{\`u}}1 {È}{{\`E}}1 {À}{{\`A}}1
}
\lstset{style=c}
\newtcblisting{codice}{listing only, nota, left=6mm,
  listing options={style=c}}
\newtcblisting{comando}{listing only, nota,
  listing options={basicstyle=\ttfamily\small, numbers=none, language={},
    columns=flexible, keepspaces=true, breaklines=true}}
\newtcblisting{risultato}{listing only, enhanced jigsaw, breakable,
  colback=black!5, colframe=black, boxrule=0.5pt, sharp corners,
  listing options={basicstyle=\ttfamily\small, numbers=none, language={},
    columns=flexible, keepspaces=true, breaklines=true}}

% --- Diagrammi di flusso (simboli standard) ---
\tikzset{
  terminale/.style = {ellipse, draw, minimum width=2.5cm, minimum height=0.9cm},
  io/.style        = {trapezium, trapezium left angle=70, trapezium right angle=110,
                      draw, minimum width=2.5cm, minimum height=0.9cm},
  processo/.style  = {rectangle, draw, minimum width=2.5cm, minimum height=0.9cm},
  decisione/.style = {diamond, draw, aspect=2, minimum width=2.5cm},
  freccia/.style   = {thick, -{Stealth}}
}

\title{Gli insiemi numerici}
\author{Marco Cerbella}
\date{Lezione 1 -- 24 settembre 2026}

\begin{document}
\maketitle

\section{Gli insiemi numerici}

Si parte da $\mathbb{N}$ e, ogni volta che un'operazione non è sempre
possibile, si allarga l'insieme:
\[
\mathbb{N} \subset \mathbb{Z} \subset \mathbb{Q} \subset \mathbb{R}.
\]

Gli insiemi numerici $\mathbb{Q}$ e $\mathbb{R}$ hanno le seguenti proprietà:
\begin{enumerate}
    \item[(R1)] è ben definita un'operazione di \textbf{somma};
    \item[(R2)] è ben definita un'operazione di \textbf{prodotto};
    \item[(R3)] è ben definita una \textbf{relazione d'ordine}: presi due numeri qualsiasi $a, b$ è sempre possibile confrontarli con ``$\leq$'', cioè vale $a \leq b$ oppure $b \leq a$.
\end{enumerate}
L'insieme dei numeri reali ha anche un'altra proprietà (l'assioma di completezza, vedi sotto):
\begin{enumerate}
    \item[(R4)] siano $A$ e $B$ due insiemi non vuoti di numeri reali con $a \leq b$ per ogni $a \in A$ e $b \in B$; allora esiste almeno un $c \in \mathbb{R}$ tale che $a \leq c \leq b$ per ogni $a \in A$ e $b \in B$.
\end{enumerate}
In particolare $\mathbb{R}$ è \emph{ordinato} (R3) ed è \emph{denso}: presi due reali qualsiasi se ne trova sempre un altro compreso tra loro. Queste proprietà permettono di rappresentare $\mathbb{R}$ come retta orientata, con una corrispondenza biunivoca tra i punti della retta e i numeri reali.

\subsection{Numeri naturali $\mathbb{N}$}

\[
\mathbb{N} = \{0, 1, 2, 3, 4, \dots\}
\]

In $\mathbb{N}$ sono definite due operazioni \emph{interne}, somma e
prodotto: sommando o moltiplicando due naturali si ottiene sempre un
naturale. Per ogni $m, n, p \in \mathbb{N}$ valgono:
\begin{enumerate}
    \item \textbf{associativa}: $(m + n) + p = m + (n + p)$ e
          $(m \cdot n) \cdot p = m \cdot (n \cdot p)$;
    \item \textbf{commutativa}: $m + n = n + m$ e $m \cdot n = n \cdot m$;
    \item \textbf{distributiva} del prodotto sulla somma:
          $m \cdot (n + p) = m \cdot n + m \cdot p$;
    \item \textbf{esistenza dello zero e dell'unità}: $0$ è elemento neutro
          della somma ($n + 0 = n$) e $1$ è elemento neutro del prodotto
          ($n \cdot 1 = n$);
    \item \textbf{leggi di cancellazione}: se $m + p = n + p$ allora $m = n$;
          se $m \cdot p = n \cdot p$ con $p \neq 0$ allora $m = n$.
\end{enumerate}

In $\mathbb{N}$ manca l'\textbf{elemento opposto}: ad esempio non esiste
$x \in \mathbb{N}$ tale che $5 + x = 0$. Quindi la sottrazione non è
sempre possibile ($3 - 5 \notin \mathbb{N}$): per questo si passa agli
interi relativi $\mathbb{Z}$.


\subsection{Interi relativi $\mathbb{Z}$}

\[
\mathbb{Z} = \{\dots, -3, -2, -1, 0, 1, 2, 3, \dots\}, \qquad
\mathbb{N} \subset \mathbb{Z}.
\]

In $\mathbb{Z}$ vale l'\textbf{esistenza dell'elemento opposto}: per ogni
$x \in \mathbb{Z}$ esiste $-x \in \mathbb{Z}$, detto \emph{opposto} di $x$,
tale che $x + (-x) = 0$ (ad esempio $3$ e $-3$). Di conseguenza la
sottrazione è sempre possibile: $a - b = a + (-b)$.

Continuano a valere tutte le proprietà dei naturali.

In $\mathbb{Z}$ manca però l'\textbf{elemento inverso}: ad esempio non
esiste $x \in \mathbb{Z}$ tale che $2 \cdot x = 1$. Quindi la divisione
non è sempre possibile ($1 : 2 \notin \mathbb{Z}$).

\subsection{Numeri razionali $\mathbb{Q}$}

\[
\mathbb{Q} = \left\{ \frac{m}{n} \;\middle|\; m, n \in \mathbb{Z},\ n \neq 0 \right\},
\qquad \mathbb{N} \subset \mathbb{Z} \subset \mathbb{Q}.
\]
Ad esempio $-2,\ 0,\ \frac{3}{2},\ -\frac{7}{4} \in \mathbb{Q}$.
Frazioni diverse possono indicare lo stesso numero:
$\frac12 = \frac24 = \frac36$.

In $\mathbb{Q}$ vale l'\textbf{esistenza dell'elemento inverso}: per ogni
$q \in \mathbb{Q}$ con $q \neq 0$ esiste $\frac{1}{q} \in \mathbb{Q}$ tale
che $q \cdot \frac{1}{q} = 1$. Quindi la divisione è sempre possibile,
\textbf{tranne per $0$}: lo zero è l'unico numero che non ha inverso
(non esiste $x$ con $0 \cdot x = 1$).

Con le sue operazioni e l'ordinamento $\leq$, $\mathbb{Q}$ è un
\emph{campo ordinato}.

\textbf{Rappresentazione decimale.} Ogni numero razionale ha una
rappresentazione decimale \textbf{finita} (es.\ $\frac14 = 0{,}25$) oppure
\textbf{periodica} (es.\ $\frac13 = 0{,}\overline{3}$), e viceversa.
I decimali infiniti \textbf{non periodici} (es.\ $1{,}41421356\dots$)
\textbf{non} sono razionali.

\textbf{Densità.} $\mathbb{Q}$ è \emph{denso}: tra due razionali ce n'è
sempre un altro. Se $a, b \in \mathbb{Q}$ con $a < b$, allora
\[
a < \frac{a + b}{2} < b,
\]
e ripetendo si trova che tra $a$ e $b$ ci sono infiniti razionali.

\textbf{I ``buchi'' di $\mathbb{Q}$.} Nonostante sia denso, $\mathbb{Q}$
non riempie la retta. La diagonale del quadrato di lato $1$ misura, per
Pitagora, $\sqrt{1^2 + 1^2} = \sqrt{2}$: riportandola sulla retta con il
compasso si ottiene un punto ben preciso, che però non corrisponde a nessun
numero razionale.

\begin{center}
\begin{tikzpicture}[scale=2.2]
  \draw[->] (-0.6,0) -- (2.5,0);
  \foreach \x in {0,1,2}
    \draw (\x,0.04) -- (\x,-0.04) node[below] {$\x$};
  \draw (0,0) rectangle (1,1);
  \draw (0,0) -- (0,1) node[midway, left] {$1$};
  \draw[thick] (0,0) -- (1,1) node[midway, above left] {$\sqrt{2}$};
  \draw[dashed] (1,1) arc[start angle=45, end angle=0, radius={sqrt(2)}];
  \fill (1.4142,0) circle (0.035);
  \node[below] at (1.4142,-0.02) {$\sqrt{2}$};
\end{tikzpicture}
\end{center}

\begin{teorema}
$\sqrt{2} \notin \mathbb{Q}$.
\end{teorema}

\begin{proof}
Per assurdo, sia $\sqrt{2} = \frac{p}{q}$ con $p, q$ senza fattori comuni
(frazione ridotta ai minimi termini). Elevando al quadrato,
$p^2 = 2q^2$: quindi $p^2$ è pari e dunque $p$ è pari, $p = 2k$.
Sostituendo, $4k^2 = 2q^2$, cioè $q^2 = 2k^2$: quindi anche $q$ è pari.
Allora $p$ e $q$ hanno il fattore comune $2$, contro l'ipotesi: assurdo.
\end{proof}

Questo ha portato all'espansione verso i numeri reali $\mathbb{R}$.


\subsection{Numeri reali $\mathbb{R}$}

\begin{definizione}
$\mathbb{R}$ è l'insieme degli allineamenti decimali con segno:
finiti, periodici oppure infiniti non periodici.
\end{definizione}

I reali che non sono razionali si dicono \emph{irrazionali}: sono i
decimali infiniti non periodici, ad esempio $\sqrt{2}$, $\pi$, $e$.
\[
\mathbb{R} = \mathbb{Q} \cup \{\text{irrazionali}\}, \qquad
\mathbb{N} \subset \mathbb{Z} \subset \mathbb{Q} \subset \mathbb{R}.
\]

In $\mathbb{R}$ valgono tutte le proprietà di $\mathbb{Q}$ (è anch'esso un
campo ordinato). Ciò che lo distingue da $\mathbb{Q}$ è un'ultima
proprietà: l'assioma di completezza.

\begin{assioma}[di completezza, o di separazione]
Siano $A, B \subseteq \mathbb{R}$ due insiemi non vuoti tali che
\[
a \leq b \quad \text{per ogni } a \in A \text{ e per ogni } b \in B.
\]
Allora esiste almeno un $c \in \mathbb{R}$ tale che
\[
a \leq c \leq b \quad \text{per ogni } a \in A \text{ e per ogni } b \in B.
\]
Il numero $c$ si chiama \emph{elemento separatore} di $A$ e $B$.
\end{assioma}

In parole: se un insieme sta tutto ``a sinistra'' di un altro, tra i due
c'è sempre un numero reale che li separa.

\begin{center}
\begin{tikzpicture}[scale=1.5]
  \draw[->] (-0.3,0) -- (6.4,0) node[right] {$\mathbb{R}$};
  \foreach \x in {0.3, 0.9, 1.6, 2.2, 2.7}
    \fill (\x,0) circle (1.6pt);
  \foreach \x in {3.4, 3.9, 4.6, 5.3, 6.0}
    \draw[fill=white] (\x,0) circle (1.6pt);
  \draw[decorate, decoration={brace, amplitude=5pt}]
    (0.3,0.2) -- (2.7,0.2) node[midway, above=6pt] {$A$};
  \draw[decorate, decoration={brace, amplitude=5pt}]
    (3.4,0.2) -- (6.0,0.2) node[midway, above=6pt] {$B$};
  \draw[thick] (3.05,0.15) -- (3.05,-0.15) node[below] {$c$};
\end{tikzpicture}
\end{center}


\begin{esempio}[Elemento separatore]
Siano $A = \{x \in \mathbb{R} \mid x \leq 3\}$ e $B = \{x \in \mathbb{R} \mid x \geq 5\}$. Ogni elemento di $A$ è minore di ogni elemento di $B$ e un elemento separatore è un qualunque $c \in [3, 5]$: ce ne sono infiniti.

Se invece $A = \{x \in \mathbb{R} \mid x < 4\}$ e $B = \{x \in \mathbb{R} \mid x \geq 4\}$, l'unico elemento separatore è $c = 4$.
\end{esempio}

\textbf{Perché in $\mathbb{Q}$ non vale.} Consideriamo i due insiemi di
razionali
\[
A = \{q \in \mathbb{Q} \mid q > 0,\ q^2 < 2\}, \qquad
B = \{q \in \mathbb{Q} \mid q > 0,\ q^2 > 2\}.
\]
Ogni elemento di $A$ è minore di ogni elemento di $B$. Gli elementi di $A$
si avvicinano a $\sqrt{2}$ da sinistra ($1;\ 1{,}4;\ 1{,}41;\ 1{,}414;
\dots$), quelli di $B$ da destra ($2;\ 1{,}5;\ 1{,}42;\ 1{,}415; \dots$).
L'unico numero che potrebbe separarli è $\sqrt{2}$, che però non è
razionale: \textbf{in $\mathbb{Q}$ non esiste un elemento separatore}.
In $\mathbb{R}$ invece l'assioma garantisce che esiste $c$, e si dimostra
che $c^2 = 2$, cioè $c = \sqrt{2}$.

\subsubsection*{Perché l'assioma di completezza è così importante}

\begin{enumerate}
    \item \textbf{La retta non ha buchi.} Grazie alla completezza, a ogni
          numero reale corrisponde un punto della retta e a ogni punto
          della retta corrisponde un numero reale (corrispondenza
          biunivoca). Per questo si parla di \emph{retta reale}, e per lo
          stesso motivo $\mathbb{R}^2$ descrive tutti i punti del piano.
    \item \textbf{Garantisce l'esistenza di numeri.} Senza completezza
          non potremmo nemmeno dire che $\sqrt{2}$ esiste. Con essa si
          dimostra che esistono le radici dei numeri positivi, $\pi$, $e$,
          i logaritmi, e che ogni allineamento decimale infinito
          rappresenta davvero un numero.
    \item \textbf{È la base di tutta l'analisi.} Limiti, continuità,
          derivate e integrali funzionano solo perché la retta è
          ``continua''. Esempio: la funzione $f(x) = x^2 - 2$ vale
          $f(1) = -1 < 0$ e $f(2) = 2 > 0$, quindi passando da negativa a
          positiva ci aspettiamo che si annulli da qualche parte. In
          $\mathbb{R}$ lo fa, in $x = \sqrt{2}$. In $\mathbb{Q}$ invece
          non si annulla mai. Teoremi come quello degli zeri, dei valori
          intermedi e di Weierstrass sono falsi in $\mathbb{Q}$ e veri in
          $\mathbb{R}$ proprio grazie alla completezza.
    \item \textbf{Forma equivalente: l'estremo superiore.} L'assioma si
          può enunciare anche così: ogni insieme non vuoto e limitato
          superiormente ha un estremo superiore in $\mathbb{R}$. È la forma
          che si usa più spesso nei teoremi di analisi.
\end{enumerate}

\textbf{Densità e completezza non sono la stessa cosa.}
\begin{center}
\begin{tabular}{l|c|c}
                & denso & completo \\ \hline
$\mathbb{Q}$    & sì    & no       \\
$\mathbb{R}$    & sì    & sì
\end{tabular}
\end{center}
Densità: tra due numeri ce n'è sempre un altro. Completezza: non ci sono
buchi. Inoltre $\mathbb{Q}$ è denso in $\mathbb{R}$: tra due numeri reali
qualsiasi c'è sempre un razionale (e anche un irrazionale).

\subsubsection*{La retta reale}

Fissati un punto $O$ (lo $0$), un verso e un'unità di misura, ogni numero
reale corrisponde a un punto della retta e viceversa.

\begin{center}
\begin{tikzpicture}[scale=1.6]
  \draw[->] (-3.5,0) -- (3.7,0) node[right] {$\mathbb{R}$};
  \foreach \x in {-3,...,3}
    \draw (\x,0.06) -- (\x,-0.06) node[below] {$\x$};
  \foreach \x/\l in {-1.5/{-\frac{3}{2}}, 0.5/{\frac{1}{2}},
                     1.4142/{\sqrt{2}}, 3.1416/{\pi}}
    {\fill (\x,0) circle (1.4pt);
     \node[above] at (\x,0.05) {$\l$};}
\end{tikzpicture}
\end{center}

I punti $-\frac32$ e $\frac12$ sono razionali; $\sqrt{2}$ e $\pi$ sono
irrazionali, ma hanno anch'essi il loro posto sulla retta: in $\mathbb{R}$
non manca nessun punto.

\end{document}
